\documentclass[12pt]{article}
\usepackage{graphicx} % Required for inserting images
\usepackage{geometry}
\geometry{a4paper, margin=2.5cm}
\usepackage[T1]{fontenc}
\usepackage{palatino}
\usepackage{amsmath, amssymb}
\usepackage{algorithm}
\usepackage{algpseudocode}
\usepackage{array}
\usepackage{titling}

\setlength{\droptitle}{-5em}
\preauthor{\begin{center}\large\vspace{0.5em}}
\postauthor{\end{center}\vspace{-0.5em}}

\title{
    {\scshape  Part IIB Revision Notes} \\ 
    \vspace{0.5em}
    \textbf{\fontsize{18}{22}\selectfont 4M24: Computational Statistics and Machine Learning}}
\author{Qianshuo (Harry) Ye}
\date{April 2026}

\usepackage{parskip}
\begin{document}
\maketitle

\section{Monte Carlo Methods}
\textbf{Central limit theorem for Monte Carlo:} The Monte Carlo estimate follows
$$
\frac{1}{N}\sum_n f(x_n) \sim \mathcal{N}\left( \mathbb{E}[f(x)], \frac{\sigma_f^2}{N} \right) \text{ as } N \to \infty
$$
where $\sigma_f^2$ is the variance of $f(x)$.

\textit{*Proof:} Define a transform to zero mean and variance
$$
\tilde{f}(x) = \frac{f(x) - \mathbb{E}[f(x)]}{\sigma_f}
$$
Define $S_N = \frac{1}{N} \sum_{n=1}^N \tilde{f}(x_n)$ and $Z_N = \frac{S_N}{\sqrt{N}}$.
The moment generating of $Z_N$ is 
\begin{align*}
M_{Z_N}(t) &= \mathbb{E}[\exp(t Z_N)] = \mathbb{E}\left[ \exp\left(t \frac{S_N}{\sqrt{N}}\right)\right] \\
&= \mathbb{E}\left[ \exp\left(\frac{t}{\sqrt{N}} \sum_{n=1}^N \tilde{f}(x_n)\right)\right] \\
&= \prod_{n=1}^N \mathbb{E}\left[ \exp\left(\frac{t}{\sqrt{N}} \tilde{f}(x_n)\right)\right] = \left( \mathbb{E}\left[ \exp\left(\frac{t}{\sqrt{N}} \tilde{f}(x)\right)\right] \right)^N
\end{align*}

Taylor expansion of the exponential function gives
$$
\exp\left(\frac{t}{\sqrt{N}} \tilde{f}(x)\right) = 1 + \frac{t}{\sqrt{N}} \tilde{f}(x) + \frac{t^2}{2N} \tilde{f}(x)^2 + O(N^{-2})
$$
Since $\mathbb{E}[\tilde{f}(x)] = 0$ and $\mathbb{E}[\tilde{f}(x)^2] = 1$, the expectation is
$$
\mathbb{E}\left[ \exp\left(\frac{t}{\sqrt{N}} \tilde{f}(x)\right)\right] = 1 + \frac{t^2}{2N} + O(N^{-2})
$$
Hence, the MGF of $Z_N$ is
$$
M_{Z_N}(t) = \left( 1 + \frac{t^2}{2N} + O(N^{-2}) \right)^N \to \exp\left(\frac{t^2}{2}\right) \text{ as } N \to \infty
$$

which is the MGF of the standard normal. Hence $Z_N \to \mathcal{N}(0, 1)$, $S_N \to \mathcal{N}(0, \frac{1}{N})$, $\frac{1}{N} \sum_{n=1}^N f(x_n) \to \mathcal{N}(\mathbb{E}[f(x)], \frac{\sigma_f^2}{N})$.

\textbf{Importance sampling:} use another probability density $q(x)$ that satisfy $q(x) \neq 0$ when $f(x) \neq 0$ to sample $x$. The expectation can be rewritten as
\begin{align*}
    \mathbb{E}_p[f(x)] &= \int f(x) p(x) dx \\
    &= \int f(x) \frac{p(x)}{q(x)} q(x) dx \\
    &= \mathbb{E}_q \left[ f(x) \frac{p(x)}{q(x)} \right]
\end{align*}

The orginal Monte Carlo estimate $\frac{1}{N} \sum_n f(x_n)$ with each $x_n \sim p(x)$ becomes $\frac{1}{N} \sum_n f(x_n) \frac{p(x_n)}{q(x_n)}$ with each $x_n \sim q(x)$. The difference in variances is
$$
\sigma_f^2 - \sigma_{\bar{f}}^2 = \int f(x)^2 p(x) dx - \int f(x)^2 \frac{p(x)^2}{q(x)} dx
$$

The biggest decrease in variance is achived when $q(x) \propto |f(x)| p(x)$.

\textbf{Control variates:} use another function $g(x)$ with known expectation $\mathbb{E}[g(x)]$ to reduce the variance of the Monte Carlo estimate. The new estimator is
$$
\hat{E}_f^c = \hat{E}_f + c (\hat{E}_g - \mathbb{E}[g(x)])
$$
where $\hat{E}_f$ and $\hat{E}_g$ are the Monte Carlo estimates of $f(x)$ and $g(x)$, respectively.

The optimal value of $c$ is
$$
c^* = -\frac{\text{Cov}(\hat{E}_f, \hat{E}_g)}{\text{Var}(\hat{E}_g)}
$$
giving the optimal estimator variance of
$$
\text{Var}(\hat{E}_f^{c^*}) = \text{Var}(\hat{E}_f) - \frac{\text{Cov}(\hat{E}_f, \hat{E}_g)^2}{\text{Var}(\hat{E}_g)}
$$

For multiple control variates, the estimator is $\hat{E}_f^c = \hat{E}_f + \sum_i c_i (\hat{E}_{g_i} - \mathbb{E}[g_i(x)])$, and the optimal coefficients are $\mathbf{c}^* = -C^{-1} \mathbf{b}$, where $C$ is the covariance matrix of $\hat{E}_{g_i}$ and $\mathbf{b}$ is the covariance vector between $\hat{E}_f$ and $\hat{E}_{g_i}$.
The optimal estimator variance is $\text{Var}(\hat{E}_f^{c^*}) = \text{Var}(\hat{E}_f) - \mathbf{b}^\top C^{-1} \mathbf{b}$.

\section{Lebesgue Integral and Measure}
\textbf{Riemann integral:} Define two summations for a function $f$ on an interval $[a, b]$ and $\Delta x_k = [x_k, x_{k+1}]$:
\begin{align*}
    S_p(x) &= \sum_{k=1}^n M_k (x_{k+1} - x_k) \text{ with } M_k = \sup_{x \in \Delta x_k} f(x) \\
    s_p(x) &= \sum_{k=1}^n m_k (x_{k+1} - x_k) \text{ with } m_k = \inf_{x \in \Delta x_k} f(x) \\
\end{align*}
Define the limit of the summations as $n \to \infty$: $S(f) = \liminf_{n \to \infty} S_p(x), s(f) = \limsup_{n \to \infty} s_p(x)$.
If $S(f) = s(f) = A$, then $A$ is the Riemann integral $\int_a^b f(x) dx$.

Limitations:
\begin{itemize}
    \item The function must be continuous almost everywhere.
    \item The construction of the integral requires a parition of the domain.
    \item The exchange of limit and integral is not always valid.
\end{itemize}

\textbf{Convergence of functions:}
\begin{itemize}
    \item \textit{Pointwise convergence:} Given $\epsilon > 0$, there exists a natural number $N(\epsilon, x)$ such that for $n > N$, $|f_n(x) - f(x)| < \epsilon$. Does not preserve continuity.
    \item \textit{Uniform convergence:} Given $\epsilon > 0$, there exists a natural number $N(\epsilon)$ such that for $n > N$, $|f_n - f| < \epsilon, \forall x \in D$. Preserves continuity.
\end{itemize}

\textbf{Lebesgue integral:} Partition the range of the function instead of the domain. For $0 \leq f_\text{min} \leq f(x) \leq f_\text{max}$, the range is partitioned into $f_\text{min} = f_1, f_2, \dots, f_n = f_\text{max}$. For each $f_k$, there are sets of values $X_k$ such that $f_k \leq f(x) < f_{k+1}$ for $x \in X_k$, and the size (measure) of each set is $\mu(X_k)$.

The Lebesgue integral is the limiting value of the sum:
$$\int f d\mu = \lim_{n \to \infty} \sum_{k=1}^n f_k \mu(X_k) = \lim_{n \to \infty} \sum_{\max |f_k - f_{k-1}| \to 0}^n f_k \mu(X_k)$$
Advantages:
\begin{itemize}
    \item Only requires a measure of sets on the domain of integration -- more general
    \item Change of order of limit and integration only requires pointwise convergence.
    \item A Riemann integrable funcion is also Lebesgue integrable, but not vice versa.
\end{itemize}

\textbf{Measure of sets:} A set $X$ and a set $\Sigma$ generated from $X$ (contains $X$, $\emptyset$, and other subsets of $X$, closed under countable complementation, union and intersections) defines a \textit{measurable set} $\{X, \Sigma\}$. Along with a measure $\mu$ it defines a \textit{measure space} $\{X, \Sigma, \mu\}$.

Rules for measure:
\begin{itemize}
    \item If $\mu(E) \geq 0$ for all $E \in \Sigma$, all elements in $\Sigma$ are measurable.
    \item $\mu(\emptyset) = 0$.
    \item If $E_i \subseteq E_j$, then $\mu(E_i) \leq \mu(E_j)$.
    \item $\mu \left(\bigcup_{i=1}^\infty E_i\right) = \sum_{i=1}^\infty \mu(E_i)$ for disjoint sets $E_i \in \Sigma$.
    \item If $(X, \Sigma_X)$ and $(Y, \Sigma_Y)$ are both measurable, then a function $f: X \to Y$ is measurable if $f^{-1}(B) \in \Sigma_X$ for all $B \in \Sigma_Y$.
\end{itemize}

\textbf{Lebesgue measure:}
\begin{itemize}
    \item For a continuous line segment on $\mathbb{R}$, the Lebesgue measure is the length of the segment. e.g. $I = [a, b]$, $\mu(I) = b - a$.
    \item The Lebesgue measure of a single point is zero. By additivity, the Lebesgue measure for the set of rational numbers $\mu(\mathbb{Q}) = 0$.
    \item The Lebesgue measure of countable sets (any set that can be put into a one-to-one correspondence with the set of integers $\mathbb{Z}$) is zero.
\end{itemize}

\textbf{Probability measures:} The \textit{probability space} $\{\Omega, \mathcal{F}, P\}$ consists of three elements:
\begin{itemize}
    \item Sample space $\Omega$: the set of all possible outcomes.
    \item $\sigma$-algebra $\mathcal{F}$: certain subsets of $\Omega$ with outcomes to be considered.
    \item Probability measure $P$: a function that assigns a probability to each event in $\mathcal{F}$, satisfying $P(\Omega) = 1$ and countable additivity.
\end{itemize}

\textbf{Random variables:} For a probability space $\{\Omega, \mathcal{F}, P\}$ and a measurable space $\{E, \mathcal{E}\}$, there is an associated measurable function $X: \Omega \to E$ that satisfies $X^{-1}(B) \in \mathcal{F}$ for all $B \in \mathcal{E}$. 
If $E = \mathbb{R}$, then $X: \Omega \to \mathbb{R}$ is a real-valued \textit{random variable} such that $\{\omega \in \Omega : X(\omega) \leq r\} \in \mathcal{F}$ for every $r \in \mathbb{R}$.

The expectation of a random variable $X$ is defined as $\mathbb{E}[X] = \int_\Omega X(\omega) dP(\omega)$.

\textbf{Radon-Nikodym theorem:} Consider two measures $\mu$ and $\nu$ on a measurable space $\{X, \Sigma\}$. If for every $A \in \Sigma$ where $\nu(A) = 0$ also have $\mu(A) = 0$, then $\mu$ is \textit{absolutely continuous} with respect to $\nu$, denoted as $\mu \ll \nu$.

Given a probability mesure $P$, a new probability measure $Q$ can be defined if there exists a function $L(\omega)$ such that $dQ(\omega) = L(\omega) dP(\omega)$.

For any measurable set $B$, $Q(B) = \int_B L(\omega) dP(\omega)$ is the probability of $B$ under the new measure $Q$.

The expectation of a function $F(\omega)$ can be computed as
$$
\int_\Omega F(\omega) dQ(\omega) = \int_\Omega F(\omega) L(\omega) dP(\omega)
$$
where the function $L(\omega) = \frac{dQ(\omega)}{dP(\omega)}$ is the Radon-Nikodym derivative if $Q \ll P$, subject to (1) $L(\omega) \geq 0$ for all $\omega \in \Omega$ and (2) $\int_\Omega L(\omega) dP(\omega) = 1$.

For a random variable $X$ on $\mathbb{R}$, the probability of a small neighborhood $B$ around the outcome $X = x \in \mathbb{R}$ is $\int_B dP = \int_B L(x) d\mu(x) = \int_B L(x) dx$. The function $L(x) = \frac{dP(x)}{dx}$ is the probability density function of $P$ with respect to the Lebesgue measure.

\section{Markov Chain Monte Carlo}
\textbf{Probability inversion:} For data $\mathbf{x} \in \mathcal{X}$ and parameters $\boldsymbol{\theta} \in \Theta$, the posterior distribution is
$
\pi (\boldsymbol{\theta} | \mathbf{x}) = \frac{p(\mathbf{x} | \boldsymbol{\theta}) \pi_0(\boldsymbol{\theta})}{\int_\Theta p(\mathbf{x} | \boldsymbol{\theta}) \pi_0(\boldsymbol{\theta}) d\boldsymbol{\theta}}
$.
Inference requires integration with respect to the posterior
$$
I(\boldsymbol{\theta}) = \mathbb{E}_{\boldsymbol{\theta} | \mathbf{x}}[f(\boldsymbol{\theta})] = \int_\Theta f(\boldsymbol{\theta}) \pi(\boldsymbol{\theta} | \mathbf{x}) d\boldsymbol{\theta}
$$
The Monte-Carlo estimate is 
$$
\hat{I}_N(\boldsymbol{\theta}) = \frac{1}{N} \sum_{n=1}^N f(\boldsymbol{\theta}_n) \text{ with } \boldsymbol{\theta}_n \sim \pi(\boldsymbol{\theta} | \mathbf{x})
$$
The estimator is unbiased and consistent as $N \to \infty$, $\hat{I}_N(\boldsymbol{\theta}) \to I(\boldsymbol{\theta})$ and the error $\sqrt{N} (\hat{I}_N(\boldsymbol{\theta}) - I(\boldsymbol{\theta})) \to \mathcal{N}(0, \sigma_f^2)$.

\textbf{Markov chains:} stochastic transitions governed by the Markov property
$$
p(\boldsymbol{\theta}^{(i)} | \boldsymbol{\theta}^{(i-1)}, \dots, \boldsymbol{\theta}^{(1)}) = T(\boldsymbol{\theta}^{(i)} | \boldsymbol{\theta}^{(i-1)})
$$
\textit{Invariant distribution}: the Markov chain converges to the invariant distribution $\pi(\boldsymbol{\theta}^{(\infty)}) \to \pi(\boldsymbol{\theta})$ irrespective of the initial state.

\textit{Irreducible:} the Markov chain can reach any state from any other state in a finite number of steps.

\textit{Aperiodic:} the Markov chain does not get trapped in cycles.

\textit{Detailed balance (reversible):} a Markov chain satisfies detailed balance if $\pi(\boldsymbol{\theta}^{(i)}) T(\boldsymbol{\theta}^{(i-1)} | \boldsymbol{\theta}^{(i)}) = \pi(\boldsymbol{\theta}^{(i-1)}) T(\boldsymbol{\theta}^{(i)} | \boldsymbol{\theta}^{(i-1)})$.
This is a sufficient condition for an invariant distribution to exist.

\subsection{Metropolis-Hastings Algorithm}
Define the \textit{transition kernel} as
$$
\underbrace{P(x, dy)}_{\text{prob. of transition from } x \text{ to } dy} = \underbrace{p(x,y) dy}_{\text{prob. of moving to $dy$ around $y$}} + \underbrace{r(x) \delta_x(dy)}_{\text{prob. of staying at } x}
$$
If $p(x, x) = 0$ and $\int_{\mathbb{R}^d} P(x,dy) = 1$, then $r(x) = 1 - \int_{\mathbb{R}^d} p(x,y) dy$ is the probability that the chain remains at $x$.

If $p(x,y)$ satisfies detailed balance $\pi(x)p(x,y) = \pi(y)p(y,x)$, then $\pi(\cdot)$ is the invariant density of the transition kernel $P(x, \cdot)$.

\textit{*Proof:} 
\begin{align*}
\int_{\mathbb{R}^d} P(x, A)\pi(x)  dx &= \int_{\mathbb{R}^d} \left[\int_A p(x,y)dy\right] \pi(x) dx + \int_{\mathbb{R}^d} r(x) \delta_x(A) \pi(x) dx \\
&= \int_A  \left[\int_{\mathbb{R}^d} \pi(x) p(x,y) dx\right] dy + \int_A r(x) \pi(x) dx \\
&= \int_A  \left[\int_{\mathbb{R}^d} \pi(y) p(y,x) dx\right] dy + \int_A r(x) \pi(x) dx \\
&= \int_A [1 - r(y)] \pi(y) dy + \int_A r(y) \pi(y) dy \\
&= \int_A \pi(y) dy = \pi^*(A)
\end{align*}

To achieve this, the \textit{proposal density} $q(x,y)$ is balanced with an \textit{acceptance probability} that satisfies $\pi(x)q(x,y)\alpha(x,y) = \pi(y)q(y,x)\alpha(y,x)$.
Set $\alpha(y,x) = 1$ then $\alpha(x,y) = \frac{\pi(y)q(y,x)}{\pi(x)q(x,y)}$ if $\pi(x)q(x,y) \ge \pi(y)q(y,x)$; otherwise $\alpha(x,y) = 1$. This gives $p_\text{MH}(x,y) = q(x,y) \alpha(x,y)$ for $x \neq y$ with the acceptance probability
$$
\alpha(x,y) = \min \left(1, \frac{\pi(y) q(y,x)}{\pi(x) q(x,y)} \right)
$$

The transition kernel is
$$
P_\text{MH}(x, dy) = q(x,y) \alpha(x,y) dy + \left(1 - \int_{\mathbb{R}^d} q(x,y) \alpha(x,y) dy\right) \delta_x(dy)
$$

For symmetric proposal density $q(x,y) = q(y,x)$, the acceptance probability simplifies to $\alpha(x,y) = \min \left(1, \frac{\pi(y)}{\pi(x)} \right)$: move to higher density regions are always accepted.

\begin{algorithm}
\caption{Metropolis-Hastings Algorithm}
\begin{algorithmic}
\State \textbf{Initialize} $x^{(1)}$
\For{$j = 1 \to N$}
    \State Simulate $y \sim q(x^{(j)}, \cdot)$ \Comment{Propose a candidate from proposal distribution}
    \State Simulate $u \sim U(0, 1)$ \Comment{Generate a random number for acceptance decision}
    \If{$u \le \alpha(x^{(j)}, y)$} \Comment{Acceptance probability $\alpha(x, y) = \min\left(1, \frac{\pi(y)q(y, x)}{\pi(x)q(x, y)}\right)$}
        \State Set $x^{(j+1)} = y$ \Comment{Accept the proposal and move to the new state}
    \Else
        \State Set $x^{(j+1)} = x^{(j)}$ \Comment{Reject the proposal and stay at the current state}
    \EndIf
\EndFor
\State \textbf{Return} $\{x^{(1)}, x^{(2)}, \dots, x^{(N)}\}$
\end{algorithmic}
\end{algorithm}

\subsection{Gibbs Sampling}
\textbf{Product transition kernels:} the product kernel $P_1(x_1, dy_1| x_2)P_2(x_2, dy_2 | y_1)$ has invariant density $\pi(y_1, y_2)$ if:
\begin{itemize}
\item the conditional kernel $P_1(x_1, dy_1 | x_2)$ has invariant density $\pi_{1|2}(\cdot | x_2)$ for fixed $x_2$
\item the conditional kernel $P_2(x_2, dy_2 | x_1)$ has invariant density $\pi_{2|1}(\cdot | x_1)$ for fixed $x_1$.
\end{itemize}

\textit{*Proof:}
\begin{align*}
&\int \int P_1(x_1, dy_1 | x_2) P_2(x_2, dy_2 | y_1) \pi(x_1, x_2) dx_1 dx_2 \\
&= \int P_2(x_2, dy_2 | y_1) \left[ \int P_1(x_1, dy_1 | x_2) \pi_{1|2}(x_1 | x_2) dx_1 \right] \pi_{2}(x_2) dx_2 \\
&= \int P_2(x_2, dy_2 | y_1) \pi^*_{1|2}(dy_1 | x_2) \pi_{2}(x_2) dx_2 \\
&= \int P_2(x_2, dy_2 | y_1) \pi_{2|1}(x_2 | y_1) \pi^*_1(dy_1) dx_2 \\
&= \pi^*_1(dy_1) \int P_2(x_2, dy_2 | y_1) \pi_{2|1}(x_2 | y_1) dx_2 \\
&= \pi^*_1(dy_1) \pi^*_{2|1}(dy_2 | y_1) = \pi^*(dy_1, dy_2)
\end{align*}

\textbf{Gibbs sampler:} set kernels to $P_1(x_1, dy_1 | x_2) = \pi^*_{1|2}(dy_1 | x_2)$ and $P_2(x_2, dy_2 | x_1) = \pi^*_{2|1}(dy_2 | x_1)$. Use exact conditional densities as the proposal denisites, i.e. for transition of $x_1$ with $x_2$ fixed, $q([x_1, x_2], [y_1, y_2]) = \pi(y_1 | x_2)$ and $q([y_1, y_2],[x_1, x_2]) = \pi(x_1 | y_2)$. As a result, the acceptance probability is $\alpha = 1$ for all moves.

\begin{algorithm}
\caption{Gibbs Sampler}
\begin{algorithmic}[1]
\For{$i = 1, \dots, N$}
    \For{$j = 1, \dots, d$} \Comment{Update one dimension at a time}
        \State $x^{(i)}_j \sim \pi(\cdot \mid x^{(i)}_1, \dots, x^{(i)}_{j-1}, x^{(i-1)}_{j+1}, \dots, x^{(i-1)}_d)$ \Comment{Sample from the conditionals}
    \EndFor
\EndFor
\State \textbf{return} $x^{(0)}, \dots, x^{(N)}$
\end{algorithmic}
\end{algorithm}

Limitation: small step sizes and poor mixing if the components are highly correlated.

\textbf{Mixture transition kernels:} mixing transitional kernels that have the same invariant density, i.e. $P(x, dy) = \sum_i w_i P_i(x, dy)$ with $\sum_i w_i = 1$.

Advantage: for multimodal distributions, the mixture kernel can have better mixing properties by combining local and global moves.

\subsection{Data Augmentation}
Augment a target density $\pi(\theta)$ with $\phi$ to form $\pi(\theta, \phi)$. Draw samples $\theta^{(n)}, \phi^{(n)} \sim \pi(\theta, \phi)$, then for each $\theta^{(n)}$ it follows $\int \pi(\theta^{(n)}, \phi) d\phi = \pi(\theta^{(n)})$.
\section{Vector, Metric and Banach Spaces}
\textbf{Vector space $V$:} a collection of vectors $x$ with axioms (group):
\begin{itemize}
    \item $x + y = y + x$ (commutativity)
    \item $(x + y) + z = x + (y + z)$ (associativity)
    \item There exists a zero vector such that $x + 0 = x$ for all $x$
    \item For each $x$, there exists an additive inverse such that $x + (-x) = 0$
    \item For $\alpha \in \mathbb{F}$, $\alpha x \in V$ for all $x \in V$ (scalar multiplication is closed)
    \item $\alpha (x + y) = \alpha x + \alpha y$ (distributivity)
    \item $\alpha (\beta x) = (\alpha \beta) x$ (associativity of scalar multiplication)
    \item $0 \times x = 0, 1 \times x = x$
\end{itemize}

\textbf{Inner product space:} space with an inner product $\langle \cdot, \cdot \rangle: V \times V \to \mathbb{F}$ satisfying:
\begin{itemize}
    \item $\langle x, y \rangle = \overline{\langle y, x \rangle}$
    \item $\langle \alpha x + \beta y, z \rangle = \alpha \langle x, z \rangle + \beta \langle y, z \rangle$ for $\alpha, \beta \in \mathbb{R}$
    \item $\langle x,x \rangle \geq 0$ and $\langle x,x \rangle = 0$ if and only if $x = 0$
    \item For $a, b \in \mathbb{C}^n, \langle a, b \rangle = \sum_{i=1}^n a^*_i b_i$
\end{itemize}
Inner product of functions: $\langle f, g \rangle = \int^b_a f(x) g(x) dx$ for $x \in [a,b] \in \mathbb{R}$.

\textbf{Vector norm:} the length of a vector, satisfying:
\begin{itemize}
    \item Schwarz inequality: $|\langle x, y \rangle| \leq \|x\| \|y\|$
    \item Triangle inequality: $\|x + y\| \leq \|x\| + \|y\|$
    \item Parallelogram law: $\|x + y\|^2 + \|x - y\|^2 = 2(\|x\|^2 + \|y\|^2)$
\end{itemize}

\textbf{Orthogonality:} two elements $x$ and $y$ are orthogonal if and only if $\langle x, y \rangle = 0$. e.g. $\sin(x)$ and $\cos(x)$ are orthogonal on $[-\pi, \pi]$.

\textbf{Orthonormal basis:} For a vector space $V$, a \textit{basis} is a \textit{finite} set of linearly independent, orthogonal unit vectors that can express any vector in $V$ as $x = \sum_{i=1}^n \alpha_i e_i$.

\textbf{Cauchy sequence:} a sequence of vectors $\{x_n\}$ is a Cauchy sequence if for any $\epsilon > 0$, there exists a natural number $N$ such that for all $m, n > N$, $\|x_n - x_m\| < \epsilon$.

If every Cauchy sequence in a vector space converges to a limit that is \textit{also in the vector space}, then the vector space is \textit{complete}.

\textbf{Metric vector space:} a vector space with a distance function $\rho(x,y)$ that maps two vectors to a non-negative real number, satisfying:
\begin{itemize}
    \item $\rho(x,y) = 0$ if and only if $x = y$
    \item $\rho(x,y) = \rho(y,x)$ (symmetry)
    \item $\rho(x,z) \leq \rho(x,y) + \rho(y,z)$ (triangle inequality)
\end{itemize}

\textbf{Normed space:} a metric space with distance function $\rho(x,y) = \|x - y\|$ induced by a non-negative norm $\|\cdot\|$ for all $x, y \in V$.

\textbf{Banach space:} \textit{a complete normed space}. All finite-dimensional normed spaces are Banach spaces.

The $l_p$ space is a Banach space where $\sum_{i=1}^\infty |x_i|^p < \infty$ for $p \geq 1$. The norm is the $l^p$ norm $\|x\|_p = \left( \sum_{i=1}^n |x_i|^p \right)^{1/p}$.

The $L_p$ space is a Banach space where $\int_a^b |f(x)|^p dx < \infty$, where $dx$ is the Lebesgue measure. The norm is the $L^p$ norm $\|f\|_p = \left( \int_a^b |f(x)|^p dx \right)^{1/p}$.
The $L^p$ norm can be used to measure the complexity of a machine learning model.

\section{Hilbert Spaces}
\textbf{Pre-Hilbert space:} an incomplete inner product space.

\textbf{Hilbert space:} \textit{a complete inner product space}. 

Any $f$ in a Hilbert space can be expressed as $f = \sum_{i=1}^\infty \langle f, e_i \rangle \, e_i$. The infinite set of infinite-dimensional orthonormal vectors $\{e_i\}$ can be a basis if and only if 
$$
\{e_i\} \text{ is complete} \iff \|f\|^2 = \sum_{i=1}^\infty |\langle f, e_i \rangle|^2 \quad \text{(Parseval's identity)}
$$
$l^2$ and $L^2$ spaces are both Hilbert spaces. The $L^2$ norm $\|f - g\|_2$ measures the closeness of functions in the mean but not necessarily pointwise. When $\|f - g\|_2 = 0$, there is \textit{equivalence almost everywhere} except for sets of measure zero. Hence, the $L^2$ space is a space of \textit{equivalence classes of functions}.

\textbf{Equivalence of $l^2$ and $L^2$ spaces:} a function $f \in L^2$ can be represented with a sequence $\{c_k = \langle f, \phi_k \rangle\} \in l^2$ by Fourier series
$$
f = \sum_{k=1}^\infty c_k \phi_k
$$
where $\{\phi_k\}$ is an orthonormal basis of $L^2$. 

$l^2$ and $L^2$ are isomorphic, i.e. there is a one-to-one correspondence between the elements of $l^2$ and $L^2$. $l^2$ is a coordinate system for $L^2$.

\textbf{Reproducing kernel Hilbert space (RKHS):} a Hilbert space of functions where point evaluation is continuous i.e. functions close in norm are also close pointwise.

The \textit{reproducing kernel function} $K_x = K(x, \cdot)$ reproduces function evaluation by inner product $f(x) = \langle f(\cdot), K_x(\cdot) \rangle$. The kernel function is symmetric and positive definite, and uniquely defines the RKHS.

\textbf{Function approximation in RKHS:} Given data $\{(x_i, y_i)\}_{i=1}^n$ where $y_i = f(x_i)$, the regularised empirical error
$$
\frac{1}{N}\sum_{n=1}^N |y_n - f(x_n)|^2 + \lambda \|f\|^2
$$
is minimised by approximating functions of the form
$$\hat{f}(\cdot) = \sum_{n=1}^N \alpha_n K(x_n, \cdot)$$
\textit{*Solution:}
\begin{align*}
\sum_{n=1}^N |y_n - \hat{f}(x_n)|^2 = \sum_{n=1}^N \left|y_n - \sum_{m=1}^N \alpha_m K(x_m, x_n)\right|^2 = (K\alpha - y)^\top (K\alpha - y) \\
\|\hat{f}\|^2 = \langle \hat{f}, \hat{f} \rangle = \sum_{n=1}^N \sum_{m=1}^N \alpha_n \alpha_m \langle K(x_n, \cdot), K(x_m, \cdot) \rangle = \sum_{n=1}^N \sum_{m=1}^N \alpha_n \alpha_m K(x_n, x_m) = \alpha^\top K \alpha
\end{align*}
The third equality follows from the reproducing property of the kernel $f = \langle f, k \rangle$.

Taking derivative with respect to $\alpha$ gives:
$$
\frac{\partial}{\partial \alpha} \left[ \frac{1}{N}(K\alpha - y)^\top (K\alpha - y) + \lambda \alpha^\top K \alpha \right] = \frac{2}{N} K(K\alpha - y) + 2\lambda K \alpha
$$
Setting the derivative to zero
\begin{align*}
\frac{1}{N} K(K\alpha - y) + \lambda K \alpha &= 0 \\
(K + N\lambda I) \alpha &= y \\
\alpha &= (K + N\lambda I)^{-1} y
\end{align*}
The term $\lambda \|f\|^2$ penalises the complexity of the function and prevents overfitting.

\section{Sobolev Spaces}
\textbf{Sobolev space:} a Sobolev space $W^k_p(\Omega)$ is a Banach space where functions $f \in L^p(\Omega)$ have \textit{weak derivatives} $D^s f \in L^p(\Omega)$ for $0 \leq s \leq k$. The norm is defined as
$$
\|f\|_{W^k_p} = \left( \sum_{s=0}^k \int_\Omega |D^s f(x)|^p \, dx \right)^{1/p}
$$

For $p = 2$, the Sobolev space $W^k_2(\Omega)$ is a Hilbert space with inner product
$$
\langle f, g \rangle_{W^k_2} = \sum_{s=0}^{k} \langle D^s f, D^s g \rangle = \sum_{s=0}^k \int_\Omega D^s f(x) D^s g(x) \, dx
$$
The norm and the inner product in $W^k_2(\Omega)$ encodes the smoothness of function.

\textbf{Weak derivatives:} for functions $f: [0,1] \to \mathbb{R}$, a function $h: [0,1] \to \mathbb{R}$ is a weak derivative of $f$ if for every differentiable function $g: [0,1] \to \mathbb{R}$ with $g(0) = g(1) = 0$,
$$
\int_0^1 f(x) g'(x) dx = -\int_0^1 h(x) g(x) dx
$$

\textbf{Function approximation by polynomial expansion:} By Fourier expansion, functions from the set $C_2[-\pi, \pi] = C[-\pi, \pi] \cap L^2[-\pi, \pi]$ can be presented as $f(x) = \sum_{k=0}^\infty c_k \exp(ikx)$ with $c_k = \frac{1}{2\pi} \int_{-\pi}^{\pi} f(x) \exp(-ikx) dx$. The $L^2$ norm of $f$ is $\|f\|^2 = \sum_{k=0}^\infty c_k^2$.

The approximating functions $f_n(x) = \sum^n_{k=0} c_k \exp(ikx)$ are in the Sobolev space $W^s_2$ with the norm
$\|f\|^2_{W^s_2} = \| D^s f \|^2_{L^2} = \sum_{k=0}^\infty k^{2s} c_k^2$.

*The approximation error is bounded as
\begin{align*}
\|f - f_n\|^2_{L^2} &= \sum_{k=n+1}^\infty c_k^2 = \sum_{k=n+1}^\infty \frac{1}{k^{2s}} k^{2s} c_k^2 \\
&< \frac{1}{n^{2s}} \sum_{k=n+1}^\infty k^{2s} c_k^2 <  \frac{1}{n^{2s}} \sum_{k=0}^\infty k^{2s} c_k^2 \\
&= \frac{1}{n^{2s}} \|f\|^2_{W^s_2}
\end{align*}
As $s$ increases, the approximation error decreases faster with increasing $n$.

For $d$-dimensional domains, the error bound is $$\|f - f_n\|^2_{L^2} < \frac{1}{dn^\frac{2s}{d}} \|f\|^2_{W^s_2}$$ which shows the curse of dimensionality -- convergence rate decreases as $d$ increases.

\section{Gaussian Measures in Hilbert Space}
\textbf{Lebesgue measure in $D$ dimensions:} $\mu = \prod_{d=1}^D \mu_d$.

There is no Lebesgue measure in infinite dimensional Hilbert or Banach spaces.

\textit{*Proof:}
\begin{itemize}
\item A Hilbert space $H$ has an associated basis $\{e_i\}_{i=1}^\infty$. 
\item Denote a ball of radius $r$ with center at $x$ as $B(x,r)$. 
\item For a ball $B(0,2)$, place a ball of radius $\frac{1}{2}$ at each basis vector $e_i$: $\{B(e_i, \frac{1}{2})\}_{i=1}^\infty$.
\item The balls do not intersect with each other: $B(e_i, \frac{1}{2}) \cap B(e_j, \frac{1}{2}) = \emptyset$ for $i \neq j$.
\item And the union of all balls is a subset of $B(0, 2)$: $\bigcup_{i=1}^\infty B(e_i, \frac{1}{2}) \subset B(0, 2)$.
\item Since the Lebesgue measure in infinite dimensions is \textit{translation invariant}, $\mu(B(e_i, \frac{1}{2})) = \mu(B(e_j, \frac{1}{2}))$ for all $i, j$.
\item \textit{Monotonicity} and \textit{countable additivity} gives $\sum_{i=1}^\infty \mu(B(e_i, \frac{1}{2})) \leq \mu(B(0, 2))$.
\item The Lebesgue measure $\mu(B(0,2)) < \infty$, but $\sum_{i=1}^{\infty} \mu(B(e_i, \frac{1}{2})) = \infty$, which is a contradiction. Therefore the Lebesgue measure is not well-defined in infinite dimensional Hilbert spaces.
\end{itemize}

\textbf{Gaussian measure:} The standard Gaussian measure is
$$
g(B) = \frac{1}{\sqrt{2\pi}} \int_B \exp\left(-\frac{x^2}{2}\right) dx
$$
For a product measure $\mu = \prod_{d=1}^\infty g_k$ on $\mathbb{R}^\infty$ to be well defined, $\exp\left(-\frac{1}{2}\sum_{k=1}^\infty x_k^2\right)$ is finite and non-zero. So $\sum_{k=1}^\infty x_k^2 < \infty$, which is the definition of $l_2$:
$$
x \in l_2 = \{x \in \mathbb{R}^\infty : \sum_{k=1}^\infty x_k^2 < \infty\}
$$
\textbf{Gaussian measure in infinite dimensions:} A Gaussian measure is a probability measure characterised by a mean function and a covariance operator.

Diagonal covariance operator: $S_x = (a_n x_n)^{\infty}_{n=1}$ for $x \in l^2$ subject to $\sum_{k=1}^\infty a_k < \infty$.

Full covariance operators: need to be \textit{trace class}, i.e. the trace $\text{Tr}(C)$ is finite
$$\text{Tr}(C) = \sum_{n=1}^\infty \langle e_n, C e_n \rangle < \infty$$
The inner product $\langle e_n, C e_n \rangle$ is
\begin{align*}
    \langle e_n, C e_n \rangle &= \mathbb{E}[(e_n, u)(u, e_n)] \\
    &= \mathbb{E}\left[\int_\Omega \int_\Omega e_n(x) [u(x)u(y)] e_n(y) \, dy \, dx\right] \\
    &= \mathbb{E}\left[\int_\Omega e_n(x) \left(\int_\Omega u(x)u(y) e_n(y) dy\right) dx\right] \\
    &= \int_\Omega e_n(x) \left(\int_\Omega \mathbb{E}[u(x)u(y)] e_n(y) dy\right) dx \\
    &= \int_\Omega e_n(x) \left(\int_\Omega c(x,y) e_n(y) dy\right) dx
\end{align*}
where $c(x,y) = \mathbb{E}[u(x)u(y)]$ is the covariance function and $C e_n(x) = \int_\Omega c(x,y) e_n(y) dy$.

\textbf{Gaussian measures in Hilbert space:} In a Hilbert space, two Gaussian measures are either equivalent or singular.

\textit{Equivalent (mutually absolutely continuous):} Two measures are equivalent if they both assign zero measure to the same sets.

\textit{Singular:} Two measures are singular w.r.t each other if one measure assigns zero to a set and the other measure assigns zero measure to the set complement.

For two Gaussian measures $\mathcal{N}(0, C)$ and $\mathcal{N}(v, C)$ in $l^2$, two measures are equivalent if $v \in \text{Image}(C^{1/2})$ i.e. $v = C^{1/2} u$ for some $u \in l^2$, and singular otherwise. The subspace for $v$ is the Cameron-Martin space (aka the reproducing kernel Hilbert space).

\textbf{Bayes' rule in Hilbert space}: The Radom-Nikodym derivative of the posterior probability measure with respect to the prior probability measure is given by
$$
\frac{d\mu^y}{d\mu^0} \propto \underbrace{P(y|x)}_{\text{likelihood}}
$$
In a infinite dimensional Hilbert space, the prior measure over functions is a Gaussian measure $\mu_0 = \mathcal{N}(m, C)$ -- formal definition of Gaussian Process prior.

\section{MCMC in Infinite-Dimensional Spaces}
\textbf{Metropolis-Hastings algorithm in infinite dimensions:} Define a proposal $$v = \sqrt{1-\beta^2} u + \beta\xi, \xi \sim \mathcal{N}(0, C)$$ such that $u$ and $v$ both follow $\mathcal{N}(0, C)$. 
Given a prior Gaussian measure $\mu^0 = \mathcal{N}(0, C)$ and a likelihood $P(y|v) = \exp(- \Phi(v))$,
the acceptance probability (replacing $\pi(v)q(v, u)$ with $\mu^y(dv)q(u, dv)$, where $\mu^y$ is the posterior measure) is 
$$\alpha(u,v) = \min\left(1, \frac{\mu^y(dv)q(v, du)}{\mu^y(du)q(u, dv)}\right) = \min\left(1, \frac{P(y|v)}{P(y|u)}\right) = \min\left(1, \exp(\Phi(u) - \Phi(v))\right)$$


\section{Langevin Dynamics}
\textbf{Langevin stochastic differential equation (SDE):} $$dX_t = -\nabla U(X_t) dt + \sqrt{2} dB_t = \nabla \log \pi(X_t) dt + \sqrt{2} dB_t$$
where $U(\cdot): \mathbb{R}^d \to \mathbb{R}$ is a potential function, $B_t$ is the Brownian motion and $\pi(x) \propto \exp(-U(x))$ is the stationary distribution of the SDE.

The measure of $X_t$, $\rho_t$, satisfies the Fokker-Planck equation
$$
\frac{\partial \rho_t}{\partial t} = \nabla \cdot (\rho_t \nabla \log \pi) + \Delta \rho_t
$$
with asymptotic solution $\rho_\infty = \pi$.

\textbf{Log-concave}: A probability measure with density $\pi$ is log-concave if
$$
\pi(\lambda x + (1-\lambda) y) \geq \pi(x)^\lambda \pi(y)^{1-\lambda} \text{ for } 0 < \lambda < 1
$$
A probability measure with density $\pi$ is m-strongly log-concave if
$$
\pi(\lambda x + (1-\lambda) y) \geq \pi(x)^\lambda \pi(y)^{1-\lambda} \exp\left( \frac{m}{2} \lambda (1-\lambda) \|x-y\|^2 \right) \text{ for } 0 < \lambda < 1
$$
i.e. there are no flat regions and the maximum is unique.

\textbf{Convergence of Langevin dynamics:} For an m-strongly log-concave $\pi$, the convergence rate to the stationary measure is exponentially fast
$$
W_2(\rho_t, \pi) \leq e^{-mt} \left\{\|x - x^*\| + \left(\frac{d}{m}\right)^\frac{1}{2}\right\}
$$
for $\pi_0 = \delta_x$ and $x^* = \arg\max_x \pi(x)$.

\textbf{Metropolis-adjusted Langevin algorithm (MALA):} Given a state of the Markov chain $X_k$, propose a new state $\bar{X}_{k+1}$ by first-order Euler discretisation with step size $\gamma$:
$$
\bar{X}_{k+1} = X_k + \gamma \nabla \log \pi(X_k) + \sqrt{2\gamma} W_k, W_k \sim \mathcal{N}(0, I)
$$
To correct the discretisation bias, accept the proposed move with probability
\begin{align*}
\alpha(X_k, \bar{X}_{k+1}) &= \min\left(1, \frac{\pi(\bar{X}_{k+1}) q(X_k | \bar{X}_{k+1})}{\pi(X_k) q(\bar{X}_{k+1}| X_k)}\right) \\
\text{where } q(\bar{X}_{k+1}| X_k) &\propto \exp\left(-\frac{1}{4\gamma} \|\bar{X}_{k+1} - X_k - \gamma \nabla \log \pi(X_k)\|^2\right)
\end{align*}

Limitation: the acceptance probability is exponentially low for high dimension, leading to poor mixing.

\textbf{Unadjusted Langevin algorithm (ULA):} Run the first-order Euler discretisation without the Metropolis-Hastings correction, i.e. $X_{k+1} = X_k + \gamma \nabla \log \pi(X_k) + \sqrt{2\gamma} W_k$. The stationary distribution is a biased distribution $\pi_\gamma$, with asymptotic bias of order $\mathcal{O}(\gamma^\frac{1}{2})$.

\textbf{Langevin schemes for Bayesian inference:} For a target distribution $\pi(x) \propto p(x)\prod_{i=1}^n p(y_i|x)$, where $p(x)$ is the prior and $p(y_i|x)$ is the likelihood for data vector $y_i$, samples can be generated from
$$
X_{k+1} = X_k + \gamma \left(\nabla \log p(X_k) + \sum_{i=1}^n \nabla \log p(y_i|X_k)\right) + \sqrt{2\gamma} W_k
$$

For MALA, the complexity is $\mathcal{O}(n)$ for both the gradient and the acceptance probability (can be avoided in ULA) computation. 

\textbf{Stochastic gradient Langevin dynamics (SGLD)}: An unbiased estimate of the gradient can be computed more efficiently by subsampling the data. At each iteration $k$, choose a random index set $\mathcal{I}_k \subset \{1, \dots, n\}$ and approximate the gradient
$$
\nabla \widehat{\log \pi(X_k)} = \nabla \log p(X_k) + \frac{n}{|\mathcal{I}_k|} \sum_{i \in \mathcal{I}_k} \nabla \log p(y_i|X_k)
$$
The complexity is $\mathcal{O}(|\mathcal{I}_k|)$, and convergence rate is similar to ULA.

\section{Energy-Based and Score-Based Models}
\textbf{Generative models:} approximate the data distribution $p_\text{data}$ with a model distribution $p_\theta$, by learning parameters $\theta$ from data $\mathcal{D} = \{x_i\}_{i=1}^N$ drawn from $p_\text{data}$. 

\textbf{Energy-based models (EBMs):} The probability of a state $x$ is defined as 
$$p(x) = \frac{1}{Z_\beta} \exp(-\beta E(x))$$
where $E(x)$ is the energy function, $\beta$ is the inverse temperature and $Z_\beta = \int \exp(-\beta E(x)) dx$ is the partition function.
EBMs paratermises the energy function, giving $p_\theta(x) \propto \exp(-E_\theta(x))$.

The EBMs are trained via maximum likelihood or contrastive divergence.

\textit{Maximum likelihood:} maximise the log-likelihood $\mathcal{L}(\theta) = \mathbb{E}_{x \sim p_\text{data}}[\log p_\theta(x)]$ via gradient methods, with gradient:
$$
\nabla_\theta \mathcal{L}(\theta) = - \mathbb{E}_{x \sim p_\text{data}}[-\nabla_\theta E_\theta(x)] + \underbrace{\mathbb{E}_{x \sim p_\theta}[-\nabla_\theta E_\theta(x)]}_{\nabla_\theta \log Z_\theta}
$$
Limitation: the second term requires sampling from the model distribution $p_\theta$, which is computationally expensive.

\textit{Contrastive divergence:}
\begin{enumerate}
    \item For each $x_i$, initialise a MCMC chain with target distribution $p_\theta$ at $x_i^{(0)}$.
    \item Run the MCMC chain for $k$ steps to obtain $x_i^{(k)}$.
    \item Compute the gradient estimate
    $$
    \hat{\mathcal{L}} = - \frac{1}{N} \sum_{i=1}^N \left( -\nabla_\theta E_\theta(x_i^{(0)}) + \nabla_\theta E_\theta(x_i^{(k)}) \right)
    $$
    \item Update parameters $\theta \leftarrow \theta + \eta \hat{\mathcal{L}}$.
\end{enumerate}
Limitations: bias in the gradient estimate and MCMC can be very slow.

\textbf{Score-based models:} Learn to match the score function $s_\theta(x) = \nabla_x \log p_\theta(x) = -\nabla_x E_\theta(x)$, bypassing the partition function.
Trained by score matching, which minimises the Fisher divergence between $p_\text{data}$ and $p_\theta$:
$$
J(\theta) = \frac{1}{2} \mathbb{E}_{x \sim p_\text{data}}\left[\|s_\theta(x) - \underbrace{\nabla \log p_\text{data}(x)}_{\text{intractable}}\|^2\right]
$$
which can be rewritten as
$$
J_\text{SM}(\theta) = \mathbb{E}_{x \sim p_\text{data}}\left[\underbrace{\text{Tr}(\nabla_x s_\theta(x))}_{\text{expensive for high-dim.}} + \frac{1}{2} \|s_\theta(x)\|^2\right]
$$

\textit{Sliced score matching:} approximate the trace of the Jacobian by projecting the score function onto random directions $v \sim \mathcal{N}(0, I)$
$$
J_\text{SSM}(\theta) = \mathbb{E}_{x \sim p_\text{data}, v \sim \mathcal{N}(0, I)}\left[v^\top \nabla_x s_\theta(x) v + \frac{1}{2} \|s_\theta(x)\|^2\right]
$$

In practice, one sample of $v$ is used for each data point:
$$
\hat{J}_\text{SSM}(\theta) = \frac{1}{N} \sum_{i=1}^N \left[v_i^\top \nabla_x s_\theta(x_i) v_i + \frac{1}{2} \|s_\theta(x_i)\|^2\right]
$$

\textbf{Sampling from generative models:} use Langevin dynamics
$$
dx_t = s_\theta(x) dt + \sqrt{2} dB_t
$$
which have stationary distribution $p_\theta$. In practice, discretise the SDE with step size $\gamma$:
$$
x^{(k+1)} = x^{(k)} + \gamma s_\theta(x^{(k)}) + \sqrt{2\gamma} z^{(k)}, \quad z^{(k)} \sim \mathcal{N}(0, I)
$$

\textbf{Parallel tempering (replica exchange):} For $\pi(x) = \exp(-U(x))$, define the tempered (smoothed) distribution $\pi_\beta(x) = \exp(-\beta U(x))$ with inverse temperature $\beta$. 
Run MCMC chains on multiple tempered distributions $\{\pi_{\beta_i}\}$ in parallel, and exchange samples between $i$-th and $j$-th chains with acceptance probability:
$$
\alpha = \min\left(1, \frac{\pi_{\beta_i}(x_t^j) \pi_{\beta_j}(x_t^i)}{\pi_{\beta_i}(x_t^i) \pi_{\beta_j}(x_t^j)}\right)
$$

\end{document}